% =====================================================================
%  INTERNAL TECHNICAL REPORT -- NOT FOR CIRCULATION
%
%  Fixed-point wordlength allocation for ADMM hardware accelerators.
%  Drafted 2026-09-06 as complete documentation of the work to date.
%
%  This is deliberately NOT a submission draft. It records refuted
%  hypotheses, withdrawn claims and scope limits alongside the results,
%  because the reasoning that killed a claim is worth more later than
%  the claim would have been. A submission version would cut Sections
%  8 and 9 and compress Section 7.
%
%  Build:  pdflatex admm_wordlength.tex   (twice, for references)
% =====================================================================
\documentclass[11pt,a4paper]{article}

\usepackage[margin=1in]{geometry}
\usepackage{amsmath,amssymb,amsthm}
\usepackage{booktabs}
\usepackage{multirow}
\usepackage{graphicx}
\usepackage{xcolor}
\usepackage[hidelinks]{hyperref}
\usepackage{listings}

\newtheorem{theorem}{Theorem}
\newtheorem{proposition}{Proposition}
\newtheorem{corollary}{Corollary}

% Flags for things that must not silently become claims.
\newcommand{\scoped}[1]{\textcolor{blue}{\textbf{[SCOPED: #1]}}}
\newcommand{\refuted}[1]{\textcolor{red}{\textbf{[REFUTED: #1]}}}
\newcommand{\todo}[1]{\textcolor{orange}{\textbf{[TODO: #1]}}}

\title{Fixed-Point Wordlength Allocation for ADMM Hardware Accelerators\\
       \large Internal technical report --- complete record including
       refuted hypotheses}
\author{}
\date{6 September 2026}

\begin{document}
\maketitle

\begin{abstract}
We study fixed-point wordlength allocation for an ADMM solver implemented on
FPGA fabric. The central observation is that a proximal operator annihilates
quantisation error on its degenerate set: the fraction $d$ of coordinates where
the operator's derivative vanishes contributes no error to the output, so the
effective noise scales with $(1-d)$ rather than the classical $q^2/12$ applied
uniformly. Measured against bit-exact simulation the resulting model achieves
ratios of $0.99$ (L1) and $1.00$ (Box); the classical model is wrong by
$3\text{--}4\times$ typically and by $900\times$ at $d = 0.999$. We derive a
range-feasibility bound on the penalty parameter, pre-registered and confirmed.
On a second problem, LASSO, run unchanged on the same hardware, the model
holds at $0.98$--$1.03$ with $d$ up to $0.77$ and an off-lattice threshold.
On routed Artix-7 silicon, narrowing the main datapath from Q2.16 to Q2.9
reduces LUTs by $46.6\%$, raises $F_{\max}$ by $28.8\%$, and cuts dynamic power
from 11\,mW to 7\,mW, giving $184 \to 117$\,nJ per solve. We report a negative
result: per-operator asymmetric lane allocation costs $4.3\%$ area at both main
widths with no power or frequency compensation. All hardware figures are
measured --- SAIF-based power, swept $F_{\max}$, gate-level bit-exact
validation --- not estimated. We state explicitly where the model does not
hold: it predicts ensemble-mean behaviour and under-predicts per-instance error
by up to $7.6$ bits.
\end{abstract}

\tableofcontents
\newpage

% =====================================================================
\section{Introduction and scope}

An accelerator for a proximal splitting method must choose a fixed-point format
for every signal in its datapath. Too few fraction bits and the solver returns
a wrong answer; too many and the fabric cost, critical path and switching
energy all grow. The classical treatment models every quantisation as an
independent additive noise source of power $q^2/12$ and propagates it through
the algorithm. This report argues that treatment is substantially wrong for
proximal operators, quantifies by how much, and validates the correction on
silicon.

\paragraph{What is claimed.} An error model in which the proximal operator's
degenerate set annihilates a measurable fraction of injected quantisation
noise; a pre-registered range-feasibility bound; a measured demonstration that
per-operator asymmetric allocation is net-negative in fabric; and a fully
measured hardware characterisation of seven configurations.

\paragraph{What is not claimed.} This is not the first error model for
fixed-point ADMM --- Jerez et al.\ have one. It is not a robustness bound. It
does not supply a mechanism for every empirical relation reported. Section~8
lists the hypotheses that were tested and failed; Section~9 lists the scope
limits.

% =====================================================================
\section{Related work and precise positioning}

\subsection{The primary prior art}

Jerez, Goulart, Richter, Constantinides, Kerrigan and Morari
(arXiv:1303.1090) implement fixed-point ADMM on FPGA with an explicit
round-off error analysis: a two-step error recurrence, Schur stability via
Lyapunov argument, and a converging upper bound. They state the bound
determines a priori the minimum number of bits for a given accuracy
specification. \textbf{Any claim of a first error model for fixed-point ADMM is
unavailable.}

The opening they leave is specific and is the basis of this work. In deriving
their bound they observe that for the projection step no arithmetic error is
introduced, and that the error \emph{can only be reduced by multiplication with
a diagonal matrix with entries componentwise in $[0,1]$}. They then set that
matrix to the identity to keep the bound conservative. That discarded diagonal
is exactly the degenerate-set annihilation quantified here.

\subsection{The second invitation}

Kinsman and Nicolici (TCAD 30(9), 2011) allocate bit-widths for iterative
computations by SMT, with conjugate gradient as their case study --- the
nearest published problem to ours. Their quantisation constraints, they write,
\emph{assume no correlation between quantisation error and value, when in
reality correlation may exist. If the nature of the correlation is known, it
can be captured into constraints.}

Their paper has sixteen citations in fifteen years and none extends it to
iterative optimisation solvers; the descendants went toward range and overflow
analysis rather than precision models for nonlinear operators.

\subsection{The gap, confirmed by the incumbents}

Ha and Sentieys state in both DATE 2020 and DATE 2023 that existing analytical
techniques are \emph{limited to modeling noise power metrics of Linear and
Time-Invariant (LTI) systems (with some extensions)}. A proximal operator is
not LTI. Two papers named the gap; a third confirms it was still open in 2023.

\subsection{Comparisons that must not be made naively}

Wu, Wallace, Mota et al.\ (J.\ Signal Process.\ Syst.\ 2022) report up to
$67\%$ LUT and $60\%$ power reduction for mixed-precision ADMM. Those figures
are \emph{not} comparable to ours. Their baseline is FP32 floating point, so
most of the gain is the float-to-fixed conversion itself; ours is Q2.16 fixed
narrowed to Q2.9, with the float-to-fixed gain already banked in the baseline.
Their power is estimated with the Xilinx Power Estimator; ours is measured from
post-route switching activity. Their precision axis is temporal (per iteration);
ours is spatial (per operator lane).

Usefully, their own Table~1 corroborates our negative result: for the ADMM
solver in fixed point, consistent FXP24 uses 6775 LUTs while mixed FXP20$\to$24
uses 6807 --- non-uniform precision costing area, unremarked.

% =====================================================================
\section{Error model}

\subsection{Degenerate-set annihilation}

Let $P_\theta$ be a proximal operator with parameter $\theta$, evaluated at
fractional wordlength $W$ with step $q = 2^{-W}$. Partition the input
coordinates into the degenerate set $D$, where $\partial P/\partial v = 0$, and
its complement.

\begin{theorem}[Operator error decomposition]
\label{thm:t1}
\begin{equation}
\mathbb{E}\lVert e_P \rVert^2 =
  \frac{|D^c|}{n}\left[L^2\frac{q^2}{12} + \sigma^2_{\theta,D^c}\right]
+ \frac{|D|}{n}\,\sigma^2_{\theta,D}
\end{equation}
\end{theorem}

Writing $d = |D|/n$, the classical model overestimates by roughly $1/(1-d)$.
On the degenerate set the operator's output does not depend on the input, so
injected perturbation is annihilated rather than propagated.

\paragraph{Validation.} \texttt{model/error\_decomp.py} measures the operator
in isolation on a $v$-distribution harvested from float ADMM runs.

\begin{table}[h]\centering
\begin{tabular}{lccc}
\toprule
operator & model/measured & classical/measured & $d$ range \\
\midrule
L1 soft-threshold & 0.99 & 3--4$\times$ & 0.00--0.84 \\
Box projection    & 1.00 & 3--4$\times$ & 0.00--0.86 \\
L2 ridge scaling  & \textbf{0.88} & --- & 0 \\
\bottomrule
\end{tabular}
\caption{Theorem~\ref{thm:t1} against bit-exact simulation. At $d = 0.999$ the
classical model is wrong by $900\times$.}
\end{table}

\scoped{L2 gives 0.88, a 12\% miss. L2 has $d = 0$, so the annihilation term is
inactive and the discrepancy lies entirely in the $\theta$ term --- the
$\gamma$ coefficient quantisation, which Corollary~1a states is not attenuated.
The annihilation mechanism is validated; the coefficient-quantisation
magnitude is not, to better than 12\%. Quoting ``0.99--1.00'' without naming
the operators would be selective.}

\subsection{Range feasibility}

\begin{proposition}[Penalty lower bound]
Weight representability requires $\rho \ge 2^{-(I-1)}$ for $I$ integer bits.
\end{proposition}

This was registered before measurement. With $I = 2$ it predicts $\rho \ge
0.5$. Measured: $\max|M| = 2.116$ at $\rho = 0.125$ (infeasible in Q2.$x$),
$1.117$ at $\rho = 0.5$ (feasible). The pre-registration is the point and
should be stated as such.

\subsection{Loop-gain relation}

Operator-level error does not determine solver-level error; the ADMM loop
amplifies it by a gain $G$. For the contractive L2 lane,
\begin{equation}
\log G = a \log R + c, \qquad R = \lVert (I-T)^{-1} \rVert,
\end{equation}
with $R$ the resolvent of the linearised iteration. Fitted on one set of
problem instances and tested on a \emph{disjoint} set, the error is 0.061 bits
in-sample and \textbf{0.046 bits held out}, inside even the 0.17-bit bar Li et
al.\ set for their SQNR model. Per instance, the model under-predicts by at
most $+0.66$ bits at the 99th percentile ($+0.73$ worst observed), so adding one
bit to its prediction covers 99\% of instances.

\refuted{For L1 and Box, fitted forms $\log G = a\log\frac{1}{1-\lVert M\rVert_2} +
b\,R_{\text{op}} + c$ with $R_{\text{L1}} = d$, $R_{\text{Box}} = d^2$ achieved
0.340 and 0.234 bits under leave-one-condition-out validation, but on the
\emph{same} instances. On disjoint instances they reach 1.38 and 0.73 bits.} For
L1 the ensemble mean is dominated by a single instance (68\% of the ensemble
error in the median cell) and shifts 0.59 bits between instance sets --- more
than the error the fit claimed. Box's held-out error is 4.7 times its own
instance-to-instance variability, so it is over-fitted rather than noisy. A
robust (geometric-mean) statistic does not rescue L1 (0.87 bits), and is
ill-posed for Box, where some instances have exactly zero error. For these two
lanes only qualitative behaviour is claimed --- loop gain grows with
conditioning and tracks the resolvent --- together with the zero-parameter
ordering predictor of Section~3.4. No mechanism is claimed for either; the
active-set churn hypothesis was tested and refuted (Section~8.2).

\paragraph{The accuracy criterion.} Li et al.\ hold their model to $<1$\,dB,
about 0.17 bits. The like-for-like comparison is Theorem~\ref{thm:t1}, an
analytical noise model against bit-exact simulation, and it meets that bar:
0.007, 0.000 and 0.092 bits for L1, Box and L2. For loop-level claims we use
0.5 bits, for three reasons: it is the resolution of an integer width decision;
no criterion can be tighter than the quantity's own repeatability, measured here
as 0.075 (L2), 0.156 (Box) and 0.593 (L1) bits between instance sets; and it is
applied to held-out instances only.

\subsection{Ordering of lane precision needs}

Which operator needs the most lane precision is not fixed. Continuous $W^*$
over a conditioning sweep shows the Box--L2 gap moving from $-0.33$ bits at
$\kappa(H^{\top}H)=2$ to $+2.33$ at 1000, crossing between 10 and 20.

The operator coefficient alone does not predict this: it gets the sign wrong in
6 of 8 resolvable Box--L2 cases and all 14 L1--L2 cases, missing by 1.2--1.5
bits. Multiplying by the loop resolvent $R=\lVert (I-T)^{-1}\rVert$, computed
from $M$ with no fitted parameter,
\begin{equation}
\Delta W_{ab} = \tfrac12\log_2\frac{A_a R_a}{A_b R_b},
\end{equation}
gives 0.288 bits and 8/8 signs for Box--L2 and places the crossover near 20.
L2's contraction pins its resolvent ($2.1\to3.2$) while L1 and Box grow
($2.7\to200$), so the contractive lane becomes relatively cheaper as the problem
hardens.

\scoped{L1--L2 is 0.621 bits and misses the two smallest gaps in sign. L1--Box
is under-predicted by 0.4--0.6 bits and the predictor is blind to $\kappa$; the
omitted $\kappa$-quantisation term of Corollary~1a is the likely cause and is
untested. The exponents fitted in the loop-gain model make Box--L2 \emph{worse}
(1.05 bits, 4/8): the zero-parameter resolvent generalises across problem
generators, the fitted model does not.}

% =====================================================================
\section{Fabric cost model}

\begin{theorem}[Lane area]
$A_{\text{lane}}(W) = bW + c$, with $b \approx 19.6$ Slice LUT/bit.
\end{theorem}

The metric matters, so it is stated: $b$ is fitted on Slice LUTs from
\texttt{report\_utilization}. On LUT cell counts the same fit gives
$b = 26.4$ and misses the held-out width by $-15.5\%$ rather than $-6.3\%$
(\texttt{model/lane\_area.py}).

\refuted{The originally proposed quadratic form
$A(W) = aW^2 + bW + c$ fails its held-out test.} Fitting on the three
no-cast points (lane width $=$ $F_{\text{MAIN}}$: $16 \to 269.9$,
$10 \to 159.0$, $9 \to 129.0$ LUT/lane) and predicting lane~8 (actual 122.0):

\begin{table}[h]\centering
\begin{tabular}{lccc}
\toprule
form & fit & predicts $W\!=\!8$ & error \\
\midrule
quadratic & exact (3 pts, 3 params) & 95.7 & $-21.6\%$ \\
linear    & 3.0\% max residual      & 114.3 & $-6.3\%$ \\
\bottomrule
\end{tabular}
\end{table}

The fitted $a$ is negative ($-1.646$): the measured slope \emph{falls} from
30\,LUT/bit ($9\to10$) to 18.5 ($10\to16$). A fabric-multiplier area law must
be convex.

\paragraph{Resolution.} The $166$ vs $35$ LUT measurement supporting a
quadratic law was taken on the $\gamma$ multiplier \emph{in isolation}, where it
does hold (4.74 measured against 4 predicted). At whole-lane granularity the
multiplier is a small share --- L1 and Box contain no multiplier at all --- so
casts, comparisons and saturating adds dominate and the total is linear. The
$W^2$ term is not refuted; it is undetectable at lane granularity on this
design. Both laws may be quoted provided the granularity is stated.

\paragraph{Cast cost, measured directly.} Two builds share lane width 8 but
differ in cast depth: 122.0\,LUT/lane at $F_{\text{MAIN}} = 10$ against
127.5 at $F_{\text{MAIN}} = 16$, giving $c_{\text{cast}} \approx
0.9$\,LUT/bit/lane.

% =====================================================================
\section{Profitability of asymmetric allocation --- a negative result}

\begin{theorem}[Profitability]
Asymmetric lane allocation is profitable only if area saved exceeds cast
overhead. With lane area linear, saving is $b\,\Delta W$ --- linear and small
--- so cast overhead dominates and asymmetry is net-negative.
\end{theorem}

Note this derivation now rests on the linear law, which survived its falsifier,
rather than on the rejected quadratic one.

\begin{table}[h]\centering
\begin{tabular}{lcc}
\toprule
matched pair (same $F_{\text{MAIN}}$, one lane narrowed) & LUT & $F_{\max}$ \\
\midrule
\texttt{asym\_derived} vs \texttt{uniform10} ($F_{\text{MAIN}}\!=\!16$)
  & $+4.3\%$ & $-5.5\%$ \\
\texttt{main9\_asym} vs \texttt{main9\_uniform} ($F_{\text{MAIN}}\!=\!9$)
  & $+4.3\%$ & $+1.3\%$ \\
\bottomrule
\end{tabular}
\end{table}

The area penalty is $+4.3\%$ at both widths, identical to one decimal place ---
a reproducible cost rather than two coincidentally similar numbers. Dynamic
power is identical across three measured pairs. The $F_{\max}$ effect is
inconsistent in sign and the $+1.3\%$ is one binary-search step, so no
frequency claim is made in either direction.

\textbf{Statement:} per-operator asymmetry costs $4.3\%$ area at both
main-datapath widths, buys no measurable dynamic power, and has no consistent
frequency effect.

% =====================================================================
\section{Hardware results}

\subsection{Method}

Artix-7 XC7A35T, out-of-context synthesis. Power from post-route gate-level
simulation writing SAIF, consumed by \texttt{report\_power}; every run
Confidence High with 72--75\% of nets matched, and bit-exact against the golden
model in the same run that produced the activity. $F_{\max}$ by binary search
on the clock constraint, seven full implementations per build, reported as
$1000/(\text{tightest period that met timing})$ --- never
$1000/(\text{period} - \text{slack})$.

\subsection{All seven builds}

\begin{table}[h]\centering\small
\begin{tabular}{lccrrrr}
\toprule
build & $F_{\text{MAIN}}$ & lanes & LUT & $F_{\max}$ (MHz) & dyn (mW) & total (mW)\\
\midrule
\texttt{uniform18}      & 16 & 16/16/16 & 6308 & 66.3 & 11 & 80 \\
\texttt{uniform10}      & 16 & 8/8/8    & 4439 & 67.6 & 10 & 78 \\
\texttt{asym\_derived}  & 16 & 8/7/8    & 4631 & 63.9 & 11 & 79 \\
\texttt{main10\_uniform}& 10 & 10/10/10 & 4081 & 79.7 &  7 & 76 \\
\texttt{main10\_lane8}  & 10 & 8/8/8    & 3665 & 68.9 &  7 & 75 \\
\textbf{\texttt{main9\_uniform}} & \textbf{9} & \textbf{9/9/9}
                                          & \textbf{3366} & \textbf{85.4}
                                          & \textbf{7} & \textbf{76} \\
\texttt{main9\_asym}    &  9 & 9/8/9    & 3510 & 86.5 &  7 & 76 \\
\bottomrule
\end{tabular}
\end{table}

\subsection{Headline}

\texttt{uniform18} $\to$ \texttt{main9\_uniform}: LUT $-46.6\%$, $F_{\max}$
$+28.8\%$ (bracketed $+27.6\%$ to $+30.4\%$ by the search resolution), dynamic
power $11 \to 7$\,mW, energy per solve $184 \to 117$\,nJ.

\paragraph{Report absolutes, not the percentage.} Power reports round to
1\,mW, so 7 and 11 carry $\pm 0.5$\,mW and the honest bound on the reduction is
29--44\%. Energy per solve is the better figure: dynamic energy per cycle is
$CV^2$ and therefore frequency-independent, so $-36\%$ holds regardless of the
operating clock a reader imagines.

\paragraph{Static power dominates.} Device static is 68\,mW of the 76--80\,mW
total and is identical across builds. Total on-chip power falls only $5.0\%$.
\textbf{Any power claim must say \emph{dynamic} explicitly or it is wrong by a
factor of seven.}

\subsection{Throughput}

Cycles $= 26K + 3$ exactly, across all seven builds, four main widths,
symmetric and asymmetric. Cycles per iteration are invariant to every width
parameter, so all performance gain comes from $F_{\max}$.

\scoped{Iteration counts are \emph{not} comparable across widths. With
$\varepsilon = 0$ the early exit fires when the residual quantises to exactly
zero, which happens sooner at coarser precision --- stagnation, not
convergence. \texttt{main9} needs 15/19/12 iterations against
\texttt{uniform18}'s 28/32/22; computing solutions per second from each build's
own iteration count bakes in a spurious $1.8\times$. Hold iterations fixed.}

\subsection{Two critical-path regimes}

Worst-path endpoints split the builds into two families with no overlap:

\begin{table}[h]\centering\small
\begin{tabular}{lccc}
\toprule
build & worst path endpoint & levels & $F_{\max}$ \\
\midrule
\texttt{uniform18}, \texttt{uniform10}, & \multirow{2}{*}{prox lane
  \texttt{z\_out\_reg}} & \multirow{2}{*}{21--24} & \multirow{2}{*}{63.9--68.9}\\
\texttt{asym\_derived}, \texttt{main10\_lane8} & & & \\
\midrule
\texttt{main10\_uniform}, \texttt{main9\_uniform}, & \multirow{2}{*}{
  \texttt{converged\_reg}} & \multirow{2}{*}{16--17} &
  \multirow{2}{*}{79.7--86.5}\\
\texttt{main9\_asym} & & & \\
\bottomrule
\end{tabular}
\end{table}

$F_{\max} = 1/\max(\text{prox path}, \text{convergence-test path})$. Narrowing
a lane below $F_{\text{MAIN}}$ inserts a narrowing cast into the prox path;
when that lifts it above the convergence-test path the build drops into the
slow regime. \texttt{main10\_lane8} and \texttt{main10\_uniform} isolate this to
one variable: same $F_{\text{MAIN}}$, same main datapath, different lane width,
different binding path, 11\,MHz apart.

\textbf{Design rule:} narrowing the main datapath raises $F_{\max}$ only while
the convergence-test path binds. Once a lane cast pushes the prox path above
it, further narrowing buys area and nothing else.

Within-family spread is $\approx 8\%$ and is implementation variance. The
family names the primal or dual residual interchangeably across builds; both
feed \texttt{converged\_reg} through the same comparison logic, so this is the
\emph{convergence-test} path and must not be described as a specific residual.

% =====================================================================
\section{A second problem: LASSO}
\label{sec:lasso}

Every result above comes from one problem family. LASSO,
$\min_x \tfrac12\|Ax-b\|^2+\lambda\|x\|_1$, maps onto the \emph{unchanged}
datapath: $M=(A^\top A+\rho I)^{-1}$, $q=A^\top b$, and the L1 lane with
$\kappa=\lambda/\rho$. We use $N=8$, a 2-sparse ground truth, Gaussian $A$ with
$m$ rows and 20\,dB SNR. Four predictions were registered in
\texttt{model/lasso.py} before it was run.

The test is harder than it looks. On the MIMO workload the L1 lane has
$d\approx0$, so Theorem~1's $(1-d)$ factor had been exercised for L1 only by
sweeping $\kappa$ with $\kappa$ placed on the lattice. LASSO yields
$d=0.62$--$0.77$ natively, with $\kappa$ wherever $\lambda$ puts it.

\begin{center}\small
\begin{tabular}{llll}
\toprule
ID & prediction & result & \\
\midrule
P-L1 & T1/meas $\in[0.9,1.1]$, native distribution & 0.981--1.028 (5 $\lambda$) & pass \\
P-L2 & $m{=}32$: $F^\star\le9$; $m{=}6$: $F^\star>9$ & 8; 10--11 & pass \\
P-L3 & support at $F^\star$ matches double $\ge99\%$ & 90--100\% per cell & \textbf{fail} \\
P-L4 & RTL bit-exact on LASSO vectors & 3 instances, $F_{\text{MAIN}}=16,9$ & pass \\
\bottomrule
\end{tabular}
\end{center}

\paragraph{P-L1.} At $\lambda/\lambda_{\max}=0.05,0.1,0.2,0.3,0.5$ the ratio of
Theorem~1's prediction to measured lane error is 0.988, 1.028, 1.019, 0.981,
0.993 (60 instances, $m=16$, $F=6$--$12$). The classical $q^2/12$ model is off
by 1.6--2.7$\times$. That is less than $1/(1-d)$ (2.7--4.4) because $\kappa$ is
off-lattice and its quantization adds error that Theorem~1's $\theta$-term
includes and the classical model does not.

\paragraph{P-L2.} Width need follows $\|M\|_2\to1/\rho$ as rows are removed:
$F^\star=8$ at $m=32$ ($\|M\|_2=0.74$), 9 at $m=16$ (0.87), 9--11 at $m=8$ (0.99),
10--11 at $m=6$ (1.00); no saturation in any cell. This is a trend over twelve
cells, not a fitted law, and $F^\star$ is not monotone for $m\le8$.

\paragraph{P-L3 failed.} Support agreement with double precision at $F^\star$ is
90--100\% per cell, not $\ge99\%$. \emph{Post hoc} (\texttt{lasso.py --e3}), all
11 disagreeing coefficients are $\le2$ LSB of $F^\star$ in magnitude (82\%
$\le1$ LSB): threshold ties at the resolution of the format. This characterises
the failure; it does not rescue the prediction.

\paragraph{P-L4.} Three instances ($m=32,16,8$) run bit-exact through the
unmodified RTL at $F_{\text{MAIN}}=16$ and 9. The $m=8$ instance recovers the
wrong support relative to ground truth in double precision as well; every
comparison in this section is against double-precision LASSO on the same
instance, never against $x_0$.

% =====================================================================
\section{Comparison with prior accelerators}
\label{sec:compare}

\textbf{Read the caveats before the numbers.} Every row below solves a
different problem, at a different size, on a different device, to a different
and mostly unstated accuracy. The table is therefore not a ranking, and this
paper claims no efficiency advantage over any row in it. We compute no
normalised metric (LUT-seconds, energy per variable) precisely because the
accuracy at which each work's solve time was obtained is not stated often
enough to make one meaningful. The table is generated by
\texttt{tools/comparison\_table.py} from \texttt{data/comparison\_lit.json},
where every entry carries a source URL and a record of who verified it.

\begin{center}\scriptsize\setlength{\tabcolsep}{3.5pt}
\begin{tabular}{llrrll}
\toprule
work & device / format & LUT & MHz & power (method) & per solve \\
\midrule
This work (baseline) & Artix-7 35T, 18 b (Q2.16) & 4673 & 66.3 & 11\,mW (SAIF, routed) & 835 cycles, 12.6 us \\
\textbf{This work} & Artix-7 35T, 11 b (Q2.9) & 2574 & 85.4 & 7\,mW (SAIF, routed) & 835 cycles, 9.8 us \\
\midrule
Jerez 2014 & Virtex-6, 18 fb & n/r & 400 & n/r (---) & 23.4 us \\
ADMIN 2021 & V7 690T, 18 b & 13392 & 263 & n/r (---) & 0.85 us \\
TASER 2016 & V7 690T, 14 b & 4790 & 232 & 600\,mW (tool est.) & 16 cyc \\
Wu 2022 & ZCU106, mixed 20-$>$24 b & 6807 & 403 & 307\,mW (XPE) & 0.07 ms \\
Hamadouche 2023 & ZCU106 (HLS), 24 b & 1822 & 157 & 44\,mW (post-impl. est.) & 7329 cyc \\
RSQP 2023 & Alveo U50, fp32 & 24245 & 300 & 19\,W (board meas.) & 31x vs CPU \\
Zhang 2025 & ZCU102, ap\_fixed$<$24,9$>$ + float & 147238 & 250 & 10.7\,W (unstated) & 4.87 ms \\
AccelMPC 2026 & Artix-7 100T, 32 b & n/r & 100 & n/r (post-route est.) & 3.3 us/iter-step \\
\bottomrule
\end{tabular}
\end{center}

\paragraph{What the table actually shows.} Three things, none of them a
speed claim.

\emph{Wordlength is assumed, not derived.} Of the fixed-point rows, every one
states a single wordlength as a given --- 18 bits, 14 bits, 24 bits, 32 bits
--- with no account of where it came from. Jerez et al.\ are the exception in
kind, offering a bound rather than a chosen number, and their own empirical
table stops at 18 fraction bits. This paper's 11 bits is derived from a model
and then confirmed on routed silicon.

\emph{Power is usually estimated.} Two rows report no power at all; three
report vendor or post-implementation estimates; one reports a card-level
measurement of a datacentre part; one states a figure without a method. Our
7\,mW is a post-route figure annotated with switching activity from a
gate-level simulation of the actual regression vectors, at Vivado
\emph{Confidence High}. We showed earlier that the vectorless number for the
same build is $\sim$10$\times$ too high, which is the size of error these
methodologies differ by.

\emph{Nobody reports the full set.} No prior ADMM accelerator we found reports
resources, activity-based power, a met $F_{\max}$, and a per-solve time
together for a fixed-point design. That, rather than any single figure, is
where this paper's measurement discipline differs.

\paragraph{Two rows that must be read carefully.} Wu et al.'s headline (67\%
LUT, 60\% power) is against an FP32 baseline; ours is Q2.16 fixed narrowed to
Q2.9, with the float-to-fixed win already in the baseline. And our own row is
the smallest problem on the smallest part: $N=8$ on a \$100 Artix-7 board, not
a ZCU106 or an Alveo. A reader comparing absolute LUTs across those parts is
comparing fabrics, not designs.

% =====================================================================
\section{Measurement discipline}

Three claims were withdrawn when vectorless estimates were replaced by
measurement, and they are recorded because the near-miss is instructive.

\begin{enumerate}
\item Vectorless power was $\approx 10\times$ too high (112 vs 11\,mW;
      78 vs 7\,mW). The ratio survived roughly intact by luck.
\item \refuted{``Asymmetric pruning cuts dynamic power 18.8\%''} --- measured,
      both configurations give 11\,mW. The difference was an artefact of the
      estimator.
\item $F_{\max}$ of 71.1\,MHz was $1000/(20\,\text{ns} - \text{slack})$ at a
      deliberately loose constraint. Swept properly it is 85.4\,MHz.
\end{enumerate}

\textbf{The near-miss.} Comparing a vectorless \texttt{uniform18} (112\,mW)
against a measured \texttt{main9} (7\,mW) would have produced a fabricated
$94\%$ power reduction. It was caught only because an object-count guard in the
SAIF flow printed zero. Every number in Section~6 is measured, and the closest
competing work estimates its power with a spreadsheet tool.

% =====================================================================
\section{Hypotheses tested and refuted}

Recorded because each would have read plausibly in a submission draft.

\subsection{Composition (Theorem 5 as originally planned)}
\refuted{2.509 bits mean error against a 0.5-bit criterion.} The original
theory placed contraction at the operator. The falsifier showed it is not
there; contraction lives at the loop. The resolvent saturates at 3.18 with
$g = 0.667$ against divergence to 47{,}669 at $g = 1$. This negative result
produced the loop-gain reframing and was predicted in advance as the claim most
likely to fail.

\subsection{Active-set churn}
The L1/Box coefficient sign difference invited a mechanism: L1's degenerate set
churns (entries crossing the soft threshold) while Box's freezes (entries
pinned to a clip bound). Churn was measured directly as the fraction of
coordinates changing degenerate-set membership per iteration.

\refuted{Two of three predictions fail.} L1 churn averages 0.0011; Box averages
0.0028 --- Box churns \emph{more}, and both are $\approx 0.1\%$, i.e.\ both
essentially frozen. Churn also fits worse than $d$ (LOO 0.781 against 0.340 for
L1). No mechanism is available and none is claimed.

\subsection{$F_{\text{MAIN}}$ dominance}
An earlier claim that narrowing the main datapath dominates lane width on all
three axes was carried by a confound: \texttt{uniform10} and
\texttt{main10\_uniform} differ in both parameters. The control build
\texttt{main10\_lane8} was added and \refuted{the frequency half fails} ---
holding lanes at 8, narrowing $F_{\text{MAIN}}$ from 16 to 10 buys $+1.9\%$
$F_{\max}$, while \emph{widening} lanes to match $F_{\text{MAIN}}$ buys
$+15.7\%$. On area both axes are real and lanes are better per bit
($-5.1\%$/bit against $-2.9\%$/bit). The correct account is the two-regime
mechanism of Section~6.5.

% =====================================================================
\section{Scope limits}

\subsection{The model predicts ensemble means, not instances}

\texttt{model/adversarial.py} attacks the fitted loop-gain model using Kinsman
and Nicolici's own construction --- right-hand sides aligned across the
eigenvectors of the system matrix --- plus a greedy search maximising model
error. Errors are reported signed, because the sign is the whole story:
positive means the model \emph{under}-predicts error, so wordlengths chosen
from it would be too few.

\begin{table}[h]\centering
\begin{tabular}{lcc}
\toprule
operator & worst under-prediction, random rhs & adversarial \\
\midrule
L1  & $+7.58$ bits & $+11.43$ bits \\
Box & $+2.16$ bits & $+3.43$ bits \\
L2  & $+0.60$ bits & $+2.25$ bits \\
\bottomrule
\end{tabular}
\end{table}

\scoped{Random instances already break the 0.5-bit criterion.} The per-instance
distribution on held-out instances (signed; $+$ means under-prediction):

\begin{center}\small
\begin{tabular}{lrrrrr}
\toprule
operator & p50 & p90 & p95 & p99 & max \\
\midrule
L2  & $-0.09$ & $+0.41$ & $+0.53$ & $+0.66$ & $+0.73$ \\
Box & $-0.92$ & $+1.67$ & $+2.76$ & $+3.46$ & $+4.05$ \\
L1  & $-3.06$ & $+0.60$ & $+1.13$ & $+1.90$ & $+4.22$ \\
\bottomrule
\end{tabular}
\end{center}
\scoped{Only L2 yields a design rule (one extra bit). The L1 and Box quantiles
are margins on models that do not generalise, so they are not one.}

\textbf{What this does not invalidate.} The Q2.9 design point was chosen with
model guidance and then verified bit-exact against the golden model on routed
silicon. The model proposes; measurement disposes. That division of labour is
the methodology and should be stated, not hidden.

\textbf{Why not a bound instead.} Converting Theorem~\ref{thm:t1} into a
guaranteed bound requires the Lyapunov or interval machinery that Jerez et al.\
and Kinsman and Nicolici already possess. Their bounds are conservative by
construction --- hundreds of bits for conjugate gradient, reducing to roughly
single-to-double precision only after heavy input-space restriction --- which
is the deficiency a predictive model exists to address. Competing there means
losing on their ground with their tools. For L1 and Box, per-instance safety therefore comes from
simulation --- which is how Q2.9 was chosen and verified --- or from the prior
art's bounds at their cost in conservatism.

\subsection{Other limits}
\begin{itemize}
\item Two LUT metrics appear in this report and must not be mixed. The
      $-46.6\%$ headline counts LUT \emph{cells}
      (\texttt{get\_cells REF\_NAME =\textasciitilde{} LUT*}): 6308 $\to$ 3366.
      Section~\ref{sec:compare} quotes Slice LUTs from
      \texttt{report\_utilization}, which is what other papers report and what
      LUT combining makes smaller: 4673 $\to$ 2574, $-44.9\%$. Both come from
      the same routed checkpoints; the ratio survives either way.
\item Power reports round to 1\,mW; this is the binding constraint on every
      power claim, including the apparent saturation at $F_{\text{MAIN}} = 10$.
\item $N = 8$, one problem family, one device. No claim is made about scaling.
\item The L1--Box ordering residual (0.4--0.6 bits, and a $\kappa$ trend of
      $-2.04$ bits per unit) is unexplained. Parameter-quantisation terms and an
      outlier explanation were both tested and refuted.
\end{itemize}

% =====================================================================
\section{Reproducibility}

No number in this report exists without a script that produces it.

\begin{center}\small
\begin{tabular}{ll}
\toprule
artefact & produced by \\
\midrule
Theorem 1 validation      & \texttt{model/error\_decomp.py} \\
Loop-gain fit             & \texttt{model/loop\_gain.py}, \texttt{loop\_gain\_l1.py} \\
Churn refutation          & \texttt{model/box\_form.py} \\
Composition refutation    & \texttt{model/generalize.py} \\
Adversarial scope test    & \texttt{model/adversarial.py} \\
LASSO (P-L1--P-L3)        & \texttt{model/lasso.py} \\
LASSO vectors (P-L4)      & \texttt{model/gen\_vectors\_lasso.py} \\
Comparison table          & \texttt{tools/comparison\_table.py} \\
Routed resource counts    & \texttt{tools/util\_table.py} \\
Golden vectors            & \texttt{model/gen\_vectors\_cfg.py} \\
Builds                    & \texttt{syn/run\_ooc.tcl} \\
SAIF power                & \texttt{syn/saif\_power.bat}, \texttt{saif\_log.tcl} \\
$F_{\max}$ sweep          & \texttt{syn/fmax\_sweep.tcl} \\
Critical-path location    & \texttt{syn/critpath.tcl} \\
Literature pass           & \texttt{tools/lit\_sweep.py} \\
\bottomrule
\end{tabular}
\end{center}

% =====================================================================
\section{Outstanding}

\begin{itemize}
\item Measured wall power (blocked on instrumentation).
\end{itemize}

\end{document}
